\documentclass[11pt,a4paper]{article}

\usepackage[utf8]{inputenc}
\usepackage[T1]{fontenc}
\usepackage[margin=2.4cm]{geometry}
\usepackage{hyperref}
\usepackage{booktabs}
\usepackage{tabularx}
\usepackage{xcolor}
\usepackage{enumitem}
\usepackage{titlesec}

\hypersetup{colorlinks=true, linkcolor=blue, urlcolor=blue, citecolor=blue}
\setlist[itemize]{leftmargin=1.5em}
\titlespacing*{\section}{0pt}{1.4em}{0.4em}
\titlespacing*{\subsection}{0pt}{1em}{0.3em}

\newcommand{\code}[1]{\texttt{#1}}

\title{Perturbation Taxonomy\\[0.3em]
\large \texttt{LLMs\_for\_NEL}: Synthetic Data Quality Control for Biomedical RE}
\author{}
\date{}

\begin{document}
\maketitle
\vspace{-3em}

\noindent
Canonical list of perturbation types we apply to BioRED gold relations to produce training data for the quality-filter classifier. The names below are the exact strings used in the \code{perturbation} column of every built CSV; treat them as a stable vocabulary.

\section*{Scope: which relation types we perturb}

The dataset focuses on the three relation types that actually have enough examples to evaluate: \code{Association}, \code{Positive\_Correlation}, \code{Negative\_Correlation}. The other five (\code{Comparison}, \code{Cotreatment}, \code{Drug\_Interaction}, \code{Bind}, \code{Conversion}) are dropped from the perturbed dataset; the CLI flag is \code{--rare-classes} (off by default; opt in to include them as an ablation). The motivation and per-class counts are in the \emph{Class-imbalance versions} section below.

\section*{How to read the table}

\begin{itemize}
\item \textbf{Family} groups perturbations that share a mechanism. Useful when slicing metrics later.
\item \textbf{Mutates} says what part of the gold tuple \code{(entity\_A, relation\_type, entity\_B)} is replaced.
\item \textbf{Why} is the LLM error mode we hope this perturbation teaches the filter to catch.
\item \textbf{Label} is the binary classifier target. \code{1} for the unmodified gold, \code{0} for everything else.
\item \textbf{Applies to} says which relation types the perturbation is meaningful for. When a perturbation cannot be meaningfully applied to a gold relation, no sample is emitted for that (gold, perturbation) pair.
\end{itemize}

\section*{Catalog}

\renewcommand{\arraystretch}{1.25}
\begin{tabularx}{\textwidth}{@{}l l X X c l@{}}
\toprule
\textbf{Name} & \textbf{Family} & \textbf{Mutates} & \textbf{Why} & \textbf{Lbl} & \textbf{Applies to} \\
\midrule
\code{gold} & positive & nothing & establishes the positive class & 1 & all kept \\
\code{label\_flip} & label & relation type, swapped to a different one (uniform random over the other kept types) & LLM picks the wrong relation type from a plausible menu & 0 & all kept \\
\code{direction\_swap} & direction & swap \code{entity\_A} and \code{entity\_B}, keep relation type & LLM reverses causality (e.g.\ ``gene up-regulates compound'' vs ``compound up-regulates gene'') & 0 & \code{Pos/Neg\_Corr.} only \\
\code{fp\_co\_related} & false pos. & replace \code{entity\_B} with another entity that is in the abstract AND has its own annotated relation (just not with \code{entity\_A}) & hardest false-positive case: both entities are independently real in the text, the model has to decide \emph{this specific pairing} is fabricated & 0 & all kept \\
\code{fp\_co\_standalone} & false pos. & replace \code{entity\_B} with an entity that is in the abstract but has no annotated relation at all & medium false-positive: the entity is co-mentioned but bears no real relation in the abstract & 0 & all kept \\
\code{fp\_external} & false pos. & replace \code{entity\_B} with an entity that is not in the abstract at all (drawn from the corpus pool) & easiest false-positive: detectable from text presence alone, weak baseline & 0 & all kept \\
\code{false\_negative} & removal & replace the relation label with \code{NoRelation} for a triple that does have a real relation & LLM drops a real relation it should have kept & 0 & all kept \\
\bottomrule
\end{tabularx}

\section*{Per-perturbation notes}

\subsection*{\code{label\_flip} applies to every kept relation type}

Frederik confirmed this on 2026-05-08: \code{label\_flip} applies to any relation type. For a gold relation of any kept type, we draw a different type uniformly from the remaining kept types. With the reduced default (\code{Association}, \code{Positive\_Correlation}, \code{Negative\_Correlation}) that means each \code{label\_flip} picks one of the other two.

\subsection*{\code{direction\_swap} applies only to \code{Positive\_Correlation} and \code{Negative\_Correlation}}

\code{Association} is \textbf{non-directional}: the relation only asserts that A and B are connected, with no claim about who acts on whom. If you swap the order of the arguments, you get back the same fact, so a ``direction-swapped Association'' sample would be labelled wrong while still being semantically correct, which is exactly the wrong training signal.

Only \code{Positive\_Correlation} and \code{Negative\_Correlation} carry direction. Those are the only types \code{direction\_swap} is applied to. Frederik made this call on 2026-05-08: Association says two things are connected but not in which direction, so direction swap only makes sense for the two correlation types.

The rare classes \code{Comparison}, \code{Cotreatment}, \code{Drug\_Interaction}, \code{Bind}, \code{Conversion} are not in the perturbed dataset by default, so the question of direction-swapping them does not arise.

\subsection*{\code{fp\_*} keep the gold relation type}

\code{fp\_co\_related}, \code{fp\_co\_standalone}, \code{fp\_external} all replace \code{entity\_B}. They \textbf{keep the gold relation's \code{relation\_type}}. We deliberately do \emph{not} randomize the relation type for fp samples, because that would overlap with \code{label\_flip} (the model could then detect fp via the wrong-relation signal rather than via the entity swap). When we keep the gold relation type, the wrong-\code{entity\_B} signal stays isolated, so per-type F1 actually measures what we claim.

\subsection*{The \code{NoRelation} $\to$ \code{real\_relation} case is \code{fp\_co\_standalone}}

On 2026-05-08 Frederik asked us to explicitly cover the case where two co-mentioned-but-unrelated entities get assigned one of the real relation labels. That case is already produced by \code{fp\_co\_standalone}, just framed from the other direction:

\begin{itemize}
\item \code{fp\_co\_standalone} starts from a real gold \code{(entity\_A, R, entity\_B)} triple,
\item replaces \code{entity\_B} with an in-abstract entity that has \textbf{no annotated relations},
\item keeps the gold's relation label \code{R}.
\end{itemize}

\noindent
The output is a sample \code{(entity\_A, R, entity\_B')} where the gold answer for that pair is \code{NoRelation}. That is exactly ``take a \code{NoRelation} pair, claim a real relation \code{R} between them''. No separate \code{no\_rel\_to\_relation} label is added; \code{fp\_co\_standalone} already covers it. Frederik just wanted it called out, not given its own label.

\subsection*{\code{false\_negative} is the inverse case}

\code{false\_negative} is the other side: take a real relation and relabel it as \code{NoRelation} (so the gold answer is \code{R} and the claimed answer is \code{NoRelation}). Together with \code{fp\_co\_standalone} this gives both directions of the \code{NoRelation} boundary.

\section*{Worked examples}

All examples below use the same paper, PMID~10491763 (first paper in BioRED Train). Title: \emph{``Hepatocyte nuclear factor-6: associations between genetic variability and type II diabetes and between genetic variability and estimates of insulin secretion.''} The relevant excerpt is:

\begin{quote}
\small
``\dots we tested the hypothesis that variability in the HNF-6 gene is associated with subsets of Type II (non-insulin-dependent) diabetes mellitus and estimates of insulin secretion in glucose tolerant subjects \dots Mutations in the coding region of the HNF-6 gene are not associated with Type II diabetes or with changes in insulin responses to glucose \dots''
\end{quote}

\noindent
The gold annotations for this paper include:
\begin{itemize}
\item \code{(HNF-6, Association, type II diabetes)}
\item \code{(glucose, Positive\_Correlation, insulin)}
\item \code{(glucose, Association, type II diabetes)}
\end{itemize}

\noindent
\textbf{Mentioned entities also include:} \code{Pro75Ala} (SequenceVariant), \code{patients} (OrganismTaxon), \code{MODY} (Gene). None of these have annotated relations in this abstract.

\subsection*{\code{gold}}
Triple: \code{(glucose, Positive\_Correlation, insulin)} \hfill Label: \code{1}\\
Supported by ``\emph{insulin secretion in glucose tolerant subjects}.''

\subsection*{\code{label\_flip}}
Triple: \code{(glucose, \textbf{Negative\_Correlation}, insulin)} \hfill Label: \code{0}\\
The abstract describes glucose-driven insulin secretion (positive). Claiming negative correlation contradicts that.

\subsection*{\code{direction\_swap}}
Triple: \code{(\textbf{insulin}, Positive\_Correlation, \textbf{glucose})} \hfill Label: \code{0}\\
The gold direction is glucose~$\to$~insulin (glucose level drives insulin). Swapped, the claim reads ``insulin drives glucose,'' which the abstract does not say. Note: this example uses \code{Positive\_Correlation}, which is one of the two relation types where \code{direction\_swap} fires.

\subsection*{\code{fp\_co\_related}}
Triple: \code{(insulin, Association, HNF-6)} \hfill Label: \code{0}\\
Both \code{insulin} and \code{HNF-6} are highlighted entities with their own real relations elsewhere in this abstract (glucose--insulin, HNF-6--diabetes), but no relation between them is annotated. The model has to know that \emph{this specific pairing} is fabricated even though both entities are independently real.

\subsection*{\code{fp\_co\_standalone}}
Triple: \code{(HNF-6, Association, Pro75Ala)} \hfill Label: \code{0}\\
\code{HNF-6}--\code{type II diabetes} is the gold pair; \code{Pro75Ala} is co-mentioned in the abstract but has \emph{no} annotated relations of its own (``\emph{Mutations in the coding region of the HNF-6 gene are not associated with Type II diabetes}''). The relation label stays as the gold's (\code{Association}). The classifier sees a triple where the gold answer for that pair is \code{NoRelation}, which is the case Frederik flagged on 2026-05-08.

\subsection*{\code{fp\_external}}
Triple: \code{(HNF-6, Association, L-thyroxine)} \hfill Label: \code{0}\\
\code{L-thyroxine} is not mentioned anywhere in this abstract (it comes from a different paper, PMID 14510914 on hypothyroidism). Detectable purely by string lookup. The relation label stays as the gold's.

\subsection*{\code{false\_negative}}
Triple: \code{(glucose, NoRelation, insulin)} \hfill Label: \code{0}\\
The gold says glucose and insulin are positively correlated. Claiming \code{NoRelation} drops that real annotation.

\section*{Per-perturbation difficulty (to be measured)}

We track macro-F1 per \code{perturbation} value, so the per-type difficulty ranking falls out of training once the filter is trained. Frederik mentioned on 2026-04-30 that in-abstract swaps will likely be harder for the model to spot than out-of-abstract ones (more semantic disambiguation needed), but we agreed to verify empirically before locking in any specific ordering.

\section*{Sampling and balance rules}

\begin{itemize}
\item \code{direction\_swap} is restricted to \code{Positive\_Correlation} and \code{Negative\_Correlation} (see the per-perturbation notes above).
\item False-positive sampling is \textbf{type-restricted by default}: only \code{(entity\_type\_a, entity\_type\_b)} pairs that appear in real BioRED relations are eligible. This avoids implausible Species/CellLine pairings the model would learn to reject trivially. The CLI flag \code{--no-type-restrict} disables it.
\item Each gold relation produces at most one sample of each perturbation type. To average out random choices (which entity got swapped, which alt-label was picked), we plan to generate \textbf{$N$ variants per gold per type} (default $N=5$; CLI: \code{--variants N}). Not yet implemented; tracked as issue~\#5.
\end{itemize}

\section*{Class-imbalance versions}

BioRED is dominated by 3 relation types. Per-class counts across the official splits:

\begin{center}
\begin{tabular}{@{}l r r r@{}}
\toprule
\textbf{Relation type} & \textbf{train} & \textbf{dev} & \textbf{test} \\
\midrule
\code{Association}          & 2192 (52.5\%) & 560 (48.2\%) & 635 (54.6\%) \\
\code{Positive\_Correlation} & 1089 (26.1\%) & 352 (30.3\%) & 325 (27.9\%) \\
\code{Negative\_Correlation} &  763 (18.3\%) & 216 (18.6\%) & 171 (14.7\%) \\
\code{Bind}                 &   61 (1.5\%)  &  19 (1.6\%)  &   9 (0.8\%)  \\
\code{Cotreatment}          &   31 (0.7\%)  &  10 (0.9\%)  &  14 (1.2\%)  \\
\code{Comparison}           &   28 (0.7\%)  &   5 (0.4\%)  &   6 (0.5\%)  \\
\code{Drug\_Interaction}    &   11 (0.3\%)  & \textbf{0}   &   2 (0.2\%)  \\
\code{Conversion}           &    3 (0.1\%)  & \textbf{0}   &   1 (0.1\%)  \\
\bottomrule
\end{tabular}
\end{center}

\noindent
The top 3 cover about 96--97\,\% of every split. \code{Drug\_Interaction} and \code{Conversion} have \textbf{zero examples in dev}, which means there is no evaluation signal for them no matter how the filter is trained. \code{Bind}, \code{Cotreatment}, \code{Comparison} are borderline at single-digit dev counts. We therefore build the dataset in two variants, both available from the same CLI:

\begin{itemize}
\item \textbf{reduced (default):} keep only \code{Association}, \code{Positive\_Correlation}, \code{Negative\_Correlation}. Built with \code{python cli.py build} (no flag) and used as the main training and evaluation set.
\item \textbf{full:} every relation type included. Built with \code{python cli.py build --rare-classes} and used as an ablation to confirm that adding the rare classes does not help (and that they are not a stealthy source of label noise).
\end{itemize}

\noindent
Both versions share the same perturbation logic; only the set of gold relations that survives into the perturbation loop changes.

\section*{Adding \code{NoRelation} as a positive-class entity pair}

\code{NoRelation} is the relation label used by the \code{false\_negative} perturbation. To make \code{NoRelation} a first-class part of the training distribution, so the filter learns that ``no relation'' is a valid ground truth and not just an absence, we also plan to sample co-mentioned entity pairs that have no annotated relation and emit them as \code{gold} with \code{relation\_type = NoRelation}, \code{label = 1}. Without this step, \code{NoRelation} only ever appears with \code{label = 0}, which teaches the model the wrong thing.

This needs conscious balancing: in real BioRED, the count of co-mentioned-but-unrelated pairs is far larger than the count of annotated relations. Cap the per-abstract count to keep \code{NoRelation} from dominating. Tracked as issue~\#7; not yet implemented.

\section*{One-thing-at-a-time spin-off datasets}

For the per-type difficulty analysis, we derive a separate spin-off CSV per perturbation type from a master sample pool. Each spin-off must differ from the others in \textbf{exactly one dimension}: the perturbation under test. Specifically:

\begin{itemize}
\item The set of gold relations covered is the same across spin-offs.
\item For each gold, the entity selection used in the perturbation (which entity got picked as the swap target, etc.) is reused across spin-offs where applicable.
\item Random seeds for variant selection are shared per gold.
\end{itemize}

\noindent
Without this discipline, observed performance differences between perturbation types are confounded by random noise in entity sampling, not by signal. Tracked as issue~\#3.

\section*{Open design questions}

\begin{itemize}
\item \textbf{Multi-ID umbrella entities} (e.g.\ \code{MODY}~$\to$~\code{3172,3651,6927}) currently get dropped (about 8\,\% of gold relations). Options: pick first ID, expand to multiple rows, skip the relation entirely. Default for now is skip; revisit before scaling up.
\item \textbf{Combo perturbations} (e.g.\ \code{label\_flip + direction\_swap} on the same gold) are not in v1. Could be added later as a harder bucket if single-perturbation accuracy saturates.
\item \textbf{\code{false\_negative} entity-swap variant:} should we also generate false-negatives where \code{entity\_B} is \emph{replaced} (so the negative is about a different pair than the gold)? Current implementation only relabels the existing pair to \code{NoRelation}.
\end{itemize}

\end{document}